\section{Related Work}
\label{sec:related}

\paragraph{Customization and identity readers.}
Personalized diffusion has matured from latent diffusion backbones deployed at
Web scale~\cite{rombach2021highresolution,shin2026latent} into a family of
subject-driven methods~\cite{cao2024controllable,wei2025personalized}.
DreamBooth binds a subject to a rare token from a handful of
images~\cite{ruiz2022dreambooth}, textual inversion encodes the subject into a
word embedding~\cite{gal2023inversion}, and multi-concept customization
composes several subjects at
once~\cite{kumari2023custom}. Zero-shot readers remove training entirely:
InstantID injects a face embedding through a decoupled cross-attention
branch~\cite{wang2024instantid}, and IP-Adapter-FaceID feeds a face identity
embedding as a prompt adapter~\cite{ye2023ipadapter}. Because these readers
run in seconds on a single photograph, a scraped portrait is immediately
actionable, which is precisely the threat model our benchmark evaluates.

\paragraph{Adversarial face and identity protection.}
Pixel-space defenses perturb images so that customization fails while the
image still looks natural. The line began with Anti-DreamBooth's
initialization-level attack on DreamBooth
fine-tuning~\cite{le2023antidreambooth} and PhotoGuard's encoder
attacks~\cite{salman2023raising}, and has since specialized along several
axes: face-specific robustness~\cite{wang2026hierarchical,qiu2026void},
face-swap pathways~\cite{wang2026safeguarding}, JPEG-robust
perturbations~\cite{fardin2026metacloakjpeg,kim2025blurguard}, real-time
identity defense~\cite{guo2024realtime}, and deepfake-oriented
shielding~\cite{jeong2025faceshield}. MetaCloak meta-learns perturbations that
survive diverse surrogate losses~\cite{liu2023metacloak}; GoodAC couples
global and local adversarial objectives for anti-customization~\cite{xu2025goodac};
GAP-Diff targets the JPEG-transcoded setting directly~\cite{zhu2025gapdiff};
IDProtector trains an adversarial noise encoder against zero-shot identity
readers~\cite{song2024idprotector}; and robust purification-resistant
objectives extend protection to diffusion-based
restoration~\cite{yang2025towards,wang2025antiaesthetics,salar2025enhancing,yahn2026personalized}.
Our pilot benchmarks the three defenses with released code or checkpoints
(GAP-Diff, MetaCloak, GoodAC) and contributes ARCP-Web as a reference defense;
the remaining variants share the same blind spot we quantify---they are
validated on pristine uploads.

\paragraph{Protection beyond perturbations.}
A parallel line changes the generation pathway rather than the pixels.
Latent-diffusion unlearning shifts denoising trajectories away from
personalized reproduction~\cite{devulapally2025latent}, Anti-Diffusion blocks
model modifications~\cite{li2025antidiffusion}, and dual-domain
anti-personalization attacks the biometric feature
space~\cite{yang2024ddap}. Immunization against instruction-driven editing and
optimization-free diffusion vaccines broaden the
attack surface~\cite{trippodo2025immunizing,ozden2024diffvax}, while
contrastive training probes the robustness of protective perturbations under
latent noise~\cite{peng2025contrastive}. Identity-unextractable synthesis,
cross-image consistency, and dual-layer defenses mix generative and
adversarial mechanisms~\cite{wang2025make,wang2025privacy,jia2025dladiff}, and
recent analyses quantify how fragile adversarial defenses remain under
distribution shift~\cite{chen2025fragile}. Extension to video protection
shows the same acquisition problem appearing in temporal
generators~\cite{huang2026delving,he2026chronolock}.

\paragraph{Unlearnable examples and data poisoning.}
Availability attacks predate customization defenses: unlearnable examples
degrade training from the data side~\cite{huang2021unlearnable}, with robust
variants~\cite{fu2022robust} and autoregressive
perturbations~\cite{sandovalsegura2022autoregressive}. Augmentation-aware
analyses then showed that benign preprocessing defeats many such
poisons~\cite{qin2023learning,jiang2023false}. The modern toolbox spans
provably unlearnable data~\cite{wang2025provably}, temporal unlearnable
videos~\cite{wu2025temporal}, trigger-free poisoning~\cite{jang2025silent},
adversarial mislabeling~\cite{wu2025feasibility}, face-privacy unlearnable
synthesis~\cite{wu2025myopia}, detection of convolution-based
poisons~\cite{li2025detecting}, augmentation-shielded
designs~\cite{gong2026armor}, and systematizations of the
field~\cite{zhang2026unlearnability,li2025survey}. Artist protection shares
the adversarial mechanics~\cite{shan2023glaze} and inherits the same
robustness questions raised by Fawkes-style cloaking~\cite{shan2020fawkes}.

\paragraph{Evaluation, recognition, and consent.}
Large-scale benchmarks standardized face-forgery evaluation under realistic
degradations~\cite{bei2024largescale,pei2024deepfake}, but no benchmark fixes
the Web acquisition channel for anti-customization. Recognition metrics
ground our measurements in face embeddings~\cite{deng2018arcface}, evaluated
on in-the-wild face collections~\cite{liu2015deep}. On the governance side,
empirical studies document how rarely Web-scraped corpora carry consent
signals~\cite{lee2025data}, and content-provenance specifications provide a
machine-readable channel for them~\cite{c2pa2026spec}. ScrapeBench connects
these threads: it fixes the acquisition contract as the unit of measurement
and pairs pixel-level robustness with consent metadata that survives when
perturbations do not.
